\documentclass[12pt]{article}

\usepackage[margin=1in]{geometry}
\usepackage[T1]{fontenc}
\usepackage{mathptmx}   % Times New Roman for text and math
\usepackage{microtype}
\usepackage{graphicx}
\usepackage{float}
\usepackage{booktabs}
\usepackage{amsmath}
\usepackage{hyperref}
\usepackage{tikz}

\usetikzlibrary{arrows.meta,positioning,shapes,fit}

\hypersetup{
  colorlinks=true,
  linkcolor=blue,
  citecolor=blue,
  urlcolor=blue
}

% ---- Cover page fields ----
\newcommand{\ReportTitle}{FireMapSim-v2 \\ An Agentic Wildfire Simulation Workspace}
\newcommand{\AuthorName}{Deepak Srinivas Govindarajan}
\newcommand{\AuthorEmail}{dg1@student.gsu.edu}
\newcommand{\CourseName}{{CSC 4980 Simulation and AI}}
\newcommand{\InstitutionName}{\textbf{Georgia State University}}
\newcommand{\SubmissionDate}{April 29, 2026}
\begin{document}

% =========================
% Cover page
% =========================
\begin{titlepage}
  \centering
  \vspace*{\fill}

  {\LARGE \textbf{\ReportTitle}\par}
  \vspace{0.25in}

  {\large \textbf{} \AuthorName \par \AuthorEmail \par}
  \vspace{0.25in}

  {\large \textbf{} \CourseName\par}
  \vspace{0.1in}
  {\large \InstitutionName\par}
  \vspace{0.2in}
  {\large \SubmissionDate\par}
  \vspace{0.2in}
  {\large \href{https://github.com/1DeepakSrinivas/firemapsim-v2}{Source Repository (GitHub)}}
  \vspace*{\fill}
\end{titlepage}

\pagenumbering{arabic}

% =========================
% Core content
% =========================

% -----------------------------------------------------------------
\section{Problem Statement}
% -----------------------------------------------------------------
Wildfire simulation tooling tends to assume a technically experienced operator who is comfortable juggling coordinate systems, weather inputs, and model-specific parameters. FireMapSim-v2 targets the gap between complex numerical models (e.g., DEVS-FIRE) and real-world decision tempo by offering:
\begin{enumerate}
    \item A map-first workspace for specifying boundaries and ignitions
    \item An AI agent that guides the scenario setup process and synchronizes decisions back into the UI.
\end{enumerate}

\subsection{Objectives}
The project aims to:
\begin{enumerate}
    \item Reduce setup friction by collecting parameters sequentially
    \item Automate data gathering from location inputs (address/ZIP/coordinates)
    \item Provide real-time visualization and replay of simulation results
    \item persist projects and results for iteration and auditability.
\end{enumerate}

\subsection{Motivation}
FireMapSim-v2 is motivated by five practical needs observed in wildfire simulation workflows:
\begin{enumerate}
    \item \textbf{An easy-to-follow workflow (vs.\ legacy FireMapSim).} Scenario setup should be a
    predictable, stepwise process rather than a collection of loosely coupled forms and parameters.
    \item \textbf{A guided, agentic flow.} An AI assistant can enforce ordering and completeness
    (location $\rightarrow$ grid $\rightarrow$ weather $\rightarrow$ ignitions $\rightarrow$ fuel
    breaks $\rightarrow$ run) while still allowing targeted manual edits when desired.
    \item \textbf{Dynamic data gathering.} Key inputs (e.g., geocoded coordinates and weather)
    should be fetched from authoritative services when the user provides minimal information such
    as a place name, ZIP code, or map click.
    \item \textbf{UI/UX-first design.} A map-first interface, immediate visual feedback, and
    clear confirmation points reduce cognitive load and help prevent configuration errors.
    \item \textbf{A Full-Stack Suite.} A unified frontend, backend, and persistence layer enables
    secure proxying to external simulators, consistent validation, reproducible runs, and saved
    project history for iteration and later review.
\end{enumerate}

\subsection{Scope}
The current scope includes a Next.js map workspace~\cite{nextjs} (Leaflet~\cite{leaflet} +
overlays), Mastra-based agentic orchestration~\cite{mastra}, weather and geospatial data
integration (Open-Meteo~\cite{openmeteo}, Nominatim~\cite{nominatim}), DEVS-FIRE proxy
endpoints~\cite{devsfire}, project persistence (PostgreSQL (Supabase)~\cite{supabase}), and authentication
(Clerk~\cite{clerk}).

% -----------------------------------------------------------------
\section{Implemented Functions}
% -----------------------------------------------------------------
FireMapSim-v2 implements the following core functions:
\begin{itemize}
  \item \textbf{Geospatial project workspace:} map visualization, project boundary definition, and
        drawing tools for ignitions and suppression/fuel-break lines.
  \item \textbf{Agentic simulation planner:} a Mastra agent (\texttt{firesim-agent}) that runs a
        strict sequential setup flow in guided mode and performs targeted edits in manual mode.
  \item \textbf{Dynamic weather integration:} server routes and tools to fetch current/hourly
        weather and populate scenario fields.
  \item \textbf{Simulation lifecycle management:} run, progress feedback, and replay (latest
        simulation overlays + metrics).
  \item \textbf{Persistence and state management:} cloud-backed project storage and a client-side
        workspace store for plan synchronization.
  \item \textbf{Authentication:} Clerk-based route protection and user identity for project
        ownership.
\end{itemize}

% Figure: Workspace entry UI — placed here as the first visual impression
\begin{figure}[H]
  \centering
  \includegraphics[width=0.95\linewidth]{images/project-report/ui-project.png}
  \caption{Workspace entry UI showing the project list and map canvas.}
  \label{fig:ui-project}
\end{figure}

% -----------------------------------------------------------------
\section{System Design}
% -----------------------------------------------------------------

\subsection{Architecture Overview}
FireMapSim-v2 is structured as a full-stack, agent-synchronized system. A user's actions live in a
shared \emph{plan} state. The agent reads a runtime context snapshot each turn and applies changes
via tools. The server mediates external calls (weather, parcels, DEVS-FIRE) and persists projects
to a PostgreSQL database (via Supabase). \\
Figure~\ref{fig:architecture} shows the major components and the primary data/control flows.

\begin{figure}[H]
  \centering
  \includegraphics[width=0.78\linewidth]{images/project-report/sys-arch.png}
  \caption{System architecture diagram showing the client workspace, Next.js API layer, Mastra
           agent, and external services.}
  \label{fig:architecture}
\end{figure}

\subsection{Component Interaction (Agent $\leftrightarrow$ UI)}
The interaction model is built around a \emph{plan-as-truth} approach:
\begin{enumerate}
  \item The client maintains a scenario plan (center coordinates, grid, weather, ignitions, fuel
        breaks).
  \item The chat host (\texttt{src/components/map/ProjectAgentChatHost.tsx}) streams messages to
        \texttt{/api/agent} and includes a runtime context object in the request body:
        \texttt{mode} (manual vs.\ guided) and \texttt{planSnapshot} (the current plan state).
  \item The Mastra agent chooses between concise manual edits or a guided sequence
        (location $\rightarrow$ grid $\rightarrow$ weather $\rightarrow$ ignitions $\rightarrow$
        fuel breaks $\rightarrow$ run).
  \item The agent API route (\texttt{src/app/api/agent/route.ts}) injects the runtime context as a
        \texttt{RUNTIME\_CONTEXT} system message and invokes the Mastra agent
        (\texttt{agentId:firesim-agent}).
  \item Confirmed values are applied via the Mastra tool \\ \texttt{update-plan}
        (\texttt{src/mastra/tools/api/updatePlan.ts}), which emits a validated patch consumed by
        the UI to update the plan deterministically.
  \item When explicitly requested, the agent issues a \texttt{run-simulation} command tool \\
        (\texttt{src/mastra/tools/api/projectCommands.ts}), which triggers the server simulation
        flow and returns replay-ready overlays and metrics.
\end{enumerate}


% Figure: Map overlays — placed after the UI interaction description
\begin{figure}[H]
  \centering
  \includegraphics[width=0.85\linewidth]{images/project-report/ui-fuel.png}
  \caption{Map overlays and scenario controls: fuel-break and ignition drawing tools active on
           the project canvas.}
  \label{fig:ui-fuel}
\end{figure}

% -----------------------------------------------------------------
\section{AI Technologies and Frameworks}
% -----------------------------------------------------------------

\subsection{Mastra}
Mastra provides the agent runtime, tools abstraction, and workflow composition~\cite{mastra}. In
FireMapSim-v2, the agent is configured with explicit operational modes (\texttt{manual}
vs.\ \texttt{chat}) and uses tool calls to mutate the shared plan state. A companion workflow
(\texttt{simulateWorkflow}) demonstrates stepwise orchestration: resolve area (including
parcel/boundary), gather weather (Open-Meteo~\cite{openmeteo}), then run or delegate to a
DEVS-FIRE execution path~\cite{devsfire}.

\subsection{OpenRouter and Qwen 3.5}
The agent uses OpenRouter as an LLM gateway~\cite{openrouter} with a Qwen~3.5 model~\cite{qwen}.
In code, the model id is configured as \texttt{openrouter/qwen/qwen3.5-397b-a17b} and is invoked
with application headers to support provider policy and analytics.

\subsection{Vercel AI SDK}
The UI uses the Vercel AI SDK for streaming chat and message state~\cite{vercelaisdk}.
Specifically, \texttt{@ai-sdk/react} provides \texttt{useChat}, and the project uses a transport
that attaches runtime context in the request body so the agent can treat the current plan as
source-of-truth each turn.

\begin{figure}[H]
  \centering
  \includegraphics[width=0.7\linewidth]{images/project-report/agent-arch.png}
  \caption{Agent architecture: \texttt{firesim-agent} tool graph and the guided sequential setup
           flow managed by Mastra.}
  \label{fig:agent-arch}
\end{figure}

% -----------------------------------------------------------------
\section{Implementation Details}
% -----------------------------------------------------------------

\subsection{Key Files and Classes}

\textbf{Agent definition.}
The \texttt{firesim-agent} is defined in
\texttt{src/mastra/agents/firesim-agent.ts}. It encodes the sequential guided flow, mode
constraints, and tool usage rules (e.g., always set both center coordinates together).

\textbf{Model provider.}
OpenRouter configuration and the Qwen~3.5 model id live in \\
\texttt{src/mastra/llm/openrouter.ts}. Vercel's AI-SDK allows this to be configured with various other providers (Anthropic, OpenAI, Google Gemini etc) as well.

\textbf{Workflow orchestration.}
\texttt{src/mastra/workflows/simulate.ts} composes steps to resolve area (including parcel
boundary), branch weather gathering, and run simulation, emitting progress events for UI
consumption.

\textbf{Workspace UI.}
The primary client workspace is \texttt{src/components/map/ProjectWorkspace.tsx}, which
coordinates: plan state, chat UI, map interactions, weather overrides, simulation execution, and
replay controls.

\textbf{Chat host.}
\texttt{src/components/map/ProjectAgentChatHost.tsx} wraps \texttt{useChat} and ensures each
request carries a runtime context body (mode + plan snapshot) to the agent API.

\textbf{Agent API route.}
\texttt{src/app/api/agent/route.ts} authenticates the user (Clerk), constructs a runtime context,
and injects it as a \texttt{RUNTIME\_CONTEXT} system message before delegating to Mastra for
streaming responses.

\textbf{DEVS-FIRE session handling}
To avoid exposing upstream simulation credentials to the browser, FireMapSim-v2 uses a server-side
session cookie for DEVS-FIRE. A connection handshake stores the upstream \texttt{userToken} in an
\texttt{HttpOnly} cookie (set server-side, not readable by client JavaScript) with a bounded TTL.
Subsequent proxy requests attach the cookie automatically; the server extracts and validates the
token before forwarding requests to DEVS-FIRE.

\textbf{DEVS-FIRE proxy and execution.}
\texttt{src/app/api/devs-fire/route.ts} provides an authenticated server-side proxy to DEVS-FIRE
endpoints (never calling the upstream directly from the browser). The execution layer
\texttt{src/lib/runDevsFireFromPlan.ts} converts the plan into DEVS-FIRE calls (including
time-step conversion, ignition dispatch, and retry handling) and returns operation logs and
metrics for replay.

\subsection{Development Environment and Setup}

\textbf{Runtime and tooling.} FireMapSim-v2 is a Bun~\cite{bun} + Next.js~\cite{nextjs} application.
Local development uses Bun as the package manager and task runner.

\textbf{Install and run.} Install dependencies and start the dev server with:
\texttt{bun install} and \texttt{bun dev}.

\textbf{Source repository.} \url{https://github.com/1DeepakSrinivas/firemapsim-v2}.

\textbf{Configuration and credentials.} The project requires runtime 
configuration for:
\begin{enumerate}
    \item LLM access (API key of LLM provider: OpenRouter, Anthropic, Google AI, Open AI etc.)
    \item Authentication
    \item Persistence / Storage
\end{enumerate}
% -----------------------------------------------------------------
\section{Demonstration Scenario}
% -----------------------------------------------------------------
 
\textbf{Deployed application link: }
\url{https://fsm2.1deepaksrinivas.com}.\\
\textbf{Demonstration video can be found at: }\href{https://studentgsu-my.sharepoint.com/:f:/g/personal/dg1_student_gsu_edu/IgAzkP_y2Yd3RZdQWtPKkcMoAdjpINCLZv88UWaDKpnDWgM?e=P2Qxod}{Link to SharePoint Folder}

\subsection{Scenario: Clearwater, Kansas --- Guided Setup Walkthrough}
This section presents a guided setup walkthrough (\texttt{mode="chat"}) of setting up the entire simulation scenario for Clearwater, Kansas.
\subsubsection{Sequence of Interactions (Summary)}
\begin{enumerate}
  \item \textbf{Guided mode handshake (runtime context).}
  \begin{itemize}
    \item \textbf{User message (example):} ``Please guide me through setting up a simulation.''
    \item \textbf{UI request:} POST to \texttt{/api/agent} with \texttt{threadId} and
          \texttt{mode="chat"}, including \texttt{planSnapshot} (the current plan state) so the agent can
          skip already-complete fields.
    \item \textbf{Agent API:} injects \texttt{RUNTIME\_CONTEXT} and invokes Mastra
          \texttt{firesim-agent}.
  \end{itemize}

  \item \textbf{Step 1 --- Location (Clearwater, Kansas).}
  \begin{itemize}
    \item \textbf{User message:} ``Set the location to Clearwater, Kansas.''
    \item \textbf{Agent response:} confirms the resolved coordinates and requests confirmation.
    \item \textbf{Tool calls:}
      \begin{itemize}
        \item \texttt{geo-geocode-address} (resolve place name via Nominatim).
        \item \texttt{update-plan} with \texttt{proj\_center\_lat} and \texttt{proj\_center\_lng}
              (and \texttt{boundaryGeoJSON} when available).
      \end{itemize}
    \item \textbf{UI:} map recenters and project boundary appears (Figure~\ref{fig:ui-project}).
  \end{itemize}

  \item \textbf{Step 2 --- Grid parameters.}
  \begin{itemize}
    \item \textbf{Agent response:} proposes defaults for grid resolution and dimensions; user confirms.
    \item \textbf{Tool call:} \texttt{update-plan} with confirmed \texttt{cellResolution},
          \texttt{cellSpaceDimension}, and \texttt{cellSpaceDimensionLat}.
    \item \textbf{UI:} the grid meta updates and becomes the constraint system for geometry edits.
  \end{itemize}

  \item \textbf{Step 3 --- Weather (dynamic fetch).}
  \begin{itemize}
    \item \textbf{User message (example):} ``Fetch current weather.'' (or confirm displayed weather)
    \item \textbf{Tool calls:}
      \begin{itemize}
        \item \texttt{weather-fetch-weather} for the project center coordinates (Open-Meteo).
        \item \texttt{update-plan} patching \texttt{windSpeed}, \texttt{windDegree}, \texttt{temperature},
              \texttt{humidity}.
      \end{itemize}
    \item \textbf{UI:} weather panels and downstream simulation inputs update consistently.
  \end{itemize}

  \item \textbf{Step 4 --- Ignitions (grid-cell geometry).}
  \begin{itemize}
    \item \textbf{User message (example):} ``Add a square ignition centered in the grid.'' 
    \item \textbf{Agent response:} translates the intent into grid-cell segments and requests confirmation.
    \item \textbf{Tool call:} \texttt{update-plan} writing ignition segments into \texttt{team\_infos[*].details}
          using canonical \texttt{start\_x/start\_y/end\_x/end\_y} fields.
    \item \textbf{UI:} ignition overlays render on the map canvas.
  \end{itemize}

  \item \textbf{Step 5 --- Fuel breaks (suppression lines).}
  \begin{itemize}
    \item \textbf{User message (example):} ``Add perimeter fuel breaks 300\,m from the edges.''
    \item \textbf{Tool call:} \texttt{update-plan} writing suppression segments into \texttt{sup\_infos}.
    \item \textbf{UI:} fuel-break overlays render and remain editable (Figure~\ref{fig:ui-fuel}).
  \end{itemize}

  \item \textbf{Step 6 --- Run and replay.}
  \begin{itemize}
    \item \textbf{Agent response:} summarizes all fields and asks for explicit confirmation.
    \item \textbf{Tool call:} \texttt{run-simulation} (optionally including \texttt{simulationTimesteps}).
    \item \textbf{Server flow:} the UI triggers server-side execution; DEVS-FIRE calls are proxied via
          \texttt{/api/devs-fire} using a server-held session cookie for the upstream token.
    \item \textbf{UI:} replay controls, overlays, and summary metrics become available
          (see Figure~\ref{fig:agent-arch} for the agent tool graph driving the flow).
  \end{itemize}
\end{enumerate}

% -----------------------------------------------------------------
\section{Conclusions}
% -----------------------------------------------------------------
FireMapSim-v2 demonstrates a practical pattern for agentic geospatial systems: treat UI state as a
formal plan, require explicit confirmation before mutation, and use structured tool calls to
synchronize changes deterministically. By combining a modern map workspace with Mastra-powered
orchestration and a robust DEVS-FIRE integration layer, the system reduces configuration friction
and supports reproducible simulation runs.

\newpage
% =========================
% Inline bibliography
% =========================
\begin{thebibliography}{99}

\bibitem{nextjs}
  \textit{Next.js Documentation}.
  \url{https://nextjs.org/docs}.

\bibitem{mastra}
  \textit{Mastra Framework}.
  \url{https://mastra.ai/}.

\bibitem{openrouter}
  \textit{OpenRouter API}.
  \url{https://openrouter.ai/docs}.

\bibitem{qwen}
  \textit{Qwen: Large Language Model Family}.
  \url{https://github.com/QwenLM}.

\bibitem{vercelaisdk}
  \textit{Vercel AI SDK}.
  \url{https://sdk.vercel.ai/}.

\bibitem{leaflet}
  \textit{Leaflet: An Open-Source JavaScript Library for Interactive Maps}.
  \url{https://leafletjs.com/}.

\bibitem{supabase}
  \textit{Supabase Documentation}.
  \url{https://supabase.com/docs}.

\bibitem{clerk}
  \textit{Clerk Authentication}.
  \url{https://clerk.com/docs}.

\bibitem{openmeteo}
  \textit{Open-Meteo Weather API}.
  \url{https://open-meteo.com/}.

\bibitem{nominatim}
  \textit{Nominatim: OpenStreetMap Search}.
  \url{https://nominatim.org/}.

\bibitem{devsfire}
  \textit{DEVS-FIRE Research API (GSU)}.
  \url{https://sims.cs.gsu.edu/sims/research/DEVSFIRE_API.html}.

\bibitem{bun}
\textit{Bun: A fast all-in-one JavaScript runtime}
\url{https://bun.com/docs}

\end{thebibliography}

\end{document}